\section{Exact verification of the finite constructions}
\label{sec:exact-reproduction}

Reproduction distinguishes the general proof from the evaluation of
a specialization. The preceding propositions use algebraic
identities, positivity and explicit convergence hypotheses. Their
computational checks use integers and rational fractions;
the lattice development additionally preserves its exact representation in
the declared algebraic extension. The record data are outputs
of the preceding construction. Executions in this
section check their geometric and operator-theoretic transport, while keeping
the provenance of their generation separate.

The first check recovers separations from minors,
verifies the Plücker relation in the triangulations and compares
chart changes with effective deformations. An invertible change
of rows multiplies every minor by the same determinant;
the marked ratios remain equal. Independent column
rescaling preserves certain cross-ratios and, in general, changes
the reconstructed separations. Evaluating both operations on the
same data provides a negative control for the equivalence
asserted.

The second check subdivides lengths by nine, preserves the
inherited markings and verifies the averaging and Gram identities.
Determinants are computed by exact elimination and compared
with their Vandermonde products. Positivity in any
finite dimension follows from the proved formula; the finite cases
executed check the implementation of that formula, its indices
and its normalizations. Counting measures and chronological
measures are kept separate so that equality of
moments does not depend on silently changing normalization.

The third check reproduces the first eight moments through
the order-four Jacobi operator. The symmetric matrix is
represented rationally in the monic-polynomial basis;
its metric is \(\operatorname{diag}(h_0,\ldots,h_3)\).
Self-adjointness in that metric, the three-term
recurrence and agreement between the matrix resolvent and its continued
fraction at \(z=2\) are verified. The result keeps explicit the observed
constant vector and the charge \(Q\), which enter the
marked spectral fidelity theorem.

The fourth group checks symplectic changes and their
Hamiltonians. Differentiation with respect to momentum recovers
velocity; the Legendre transform is checked as a polynomial
or Laurent identity. The inertia of a quadratic
form and positivity of a quantum Hamiltonian
belong to their respective spaces. The Clifford
representation is checked through its anticommutators,
and composition with the mass operator is proved by
squaring the Dirac operator. These equalities
identify each term and its cross-products.

Finally, the preceding checks of triangulation,
canonical form, Schur reduction, mass
congruence and lattice construction are preserved. In Yang--Mills, the
propositions on the common measure, refinement and reconstruction
are set out in the text with their premises and falsifiers. A
finite matrix check does not replace passage to the limit
or field reconstruction: their domains are recorded
separately in the reproducible delivery.

The source package links each program to the statement
it evaluates and preserves its complete output. Conditions are
checked through explicit branches, so that they
remain active when the interpreter is run in
optimized mode. They are also reproduced in isolated mode, without
loading user packages. The correspondence between
construction, proof and execution can be inspected
step by step without turning a program's output into
an assertion of broader scope than what it checks.
